\section{A Conductance Obstruction to Uniformly Fast Local Mixing}
\label{sec:conductance-obstruction}

The results in this section are proved below.  They were supplied in the
November 2026 conductance draft and expanded during the accompanying internal
proof audit.  They concern the global uniform law $\pi_{N,m}$ on
$\Omega(N,m)$, with products acting rightmost first.  They refute a mixing
bound polynomial in $N,m$ for the local chain over its entire parameter
range.  They do not give an impossibility theorem for direct sampling.

For comparison with the project kernel, define the adjacent pair heat bath
\[
 P_{N,m}=\frac12 I+
 \frac1{2(m-1)}\sum_{i=1}^{m-1}H_{i,2},\qquad m\ge2,
\]
where each conditional is uniform and windows do not wrap around.
For any stationary kernel $L$, write
\[
 Q_L(A,A^c)=\sum_{x\in A}\pi_{N,m}(x)L(x,A^c),\qquad
 \Phi(L)=\min_{0<\pi_{N,m}(A)\le1/2}
       \frac{Q_L(A,A^c)}{\pi_{N,m}(A)}.
\]
Its support graph has one edge for each factor, with multiplicities.
For $m=2h$, let $Z_k$ count components containing $2k$ gates, and let
$A_K$ be the event that every component has at most $2K$ gates.
Isolated vertices have zero gates.  Each nontrivial component subword is
itself an identity word, since distinct components act on disjoint labels.

\begin{theorem}[Balanced bottleneck and bounded-slot obstruction]
\label{thm:balanced-bottleneck}
There are absolute constants $c,C>0$ such that, for every sufficiently large
even $m$, some integer $N=N_m$ satisfies
\[
 \frac{c m^{3/2}}{(\log m)^2}\le N\le C m^{3/2}.
\]
At these parameters there is a cap $K\asymp N/\sqrt m$ with
$\pi_{N,m}(A_K)=1/2+o(1)$.  For the pair kernel,
\[
 \frac{Q_P(A_K,A_K^c)}{\min(\pi_{N,m}(A_K),\pi_{N,m}(A_K^c))}
 \le \exp(-cN/\sqrt m).
\]
More generally, for each fixed integer $b\ge1$, every uniform-stationary
kernel changing at most $b$ positions per step satisfies the same bound
with a constant $c_b>0$.  For some $c'_b>0$, at time
\[
 t_m=\left\lfloor\exp\!\left(\frac{c'_b m}{(\log m)^2}\right)\right\rfloor
\]
a $1-o(1)$ fraction of individual starting words have total-variation
distance at least $1/2-o(1)$ from uniform.  In particular this applies to
$\Klocal$, with $b=5$, and refutes a universal polynomial mixing bound in
$N$ and $m$ for that chain.
\end{theorem}

The quantifier on $N$ is existential: the theorem neither locates a sharp
transition nor asserts this lower bound at every $N$ in the displayed
interval.  All logarithms here are natural.

\subsection{Connected-word counts and switching}

\begin{lemma}[Connected-map bound]
\label{lem:connected-map-bound}
Let $U_{n,k}$ count length-$2k$ identity words on $[n]$ whose used labels
form one connected component.  For an absolute constant $B$ and $k\ge1$,
\[
 U_{n,k}\le(2k)!B^k n^{k+1}
                 \sum_{g=0}^k(k^2/n)^g.                 \tag{CB.1}
\]
\end{lemma}
\begin{proof}
A word with $v$ used labels defines an oriented ribbon map: order the
incident half-edges at each vertex by gate index, cyclically.  Tracing face
boundaries gives the cycles of the product (reversing the convention
reverses the product, with the same identity conclusion).  Thus an identity
word gives $v$ faces.  This is the pure Hurwitz-map correspondence of
\cite[Definition 1.3 and equation (10)]{CLW2024}.  Euler's formula gives
$2v-2k=2-2g$, so $v=k+1-g$.

Here is an explicit upper-bound encoding, uniform in the genus.  Delete the
primal edges of a dual spanning tree, merging the faces into one disk and
leaving a connected one-face map.  Choose a spanning tree of that map.
Its rooted plane-tree shape has at most $4^{v-1}$ possibilities.  The
one-face map has $2g$ further edges.  Distribute their $4g$ ends as ordered
lists in the tree's corners, then pair the ends.  Stars and bars gives
$2^{O(k)}$ distributions; the number of pairings is
$(4g-1)!!\le(4k)^{2g}$, with value one when $g=0$.

The deleted $v-1$ dual-tree edges can now be reinserted in the disk face.
Distribute their ends in its boundary corners, in $2^{O(k)}$ ways.  The
ends admit at most $4^{v-1}$ noncrossing pairings.  This reconstructs every
original map for at least one choice of the two trees; extra encodings
only increase the bound.  Starting-corner and root choices cost at most a
polynomial in $k$, absorbed into $B^k$.  Consequently the number of
unlabelled maps is at most $B^k k^{2g}$.  Injective vertex labels and
distinct chronological edge labels contribute at most $n^v(2k)!$.
Summing over $g$ proves the bound.  Only $v\ge2$ can actually occur;
extending the sum to $g=k$ is harmless.
\end{proof}

\begin{lemma}[Conditional component switching]
\label{lem:component-switching}
There are absolute $C_0,C_1,c_0>0$ such that, if $m=2h\le c_0N$ and
$2\le k\le K\le h$, then
\[
 \mathbb E[Z_k\mid A_K]\le\Lambda_k,
 \qquad
 \Lambda_k=C_0N\left[C_1\max\left\{\frac mN,
                                      \frac{mk^2}{N^2}\right\}\right]^k.
                                                        \tag{CB.2}
\]
Moreover, writing $D_{N,h}$ for the number of words consisting of $h$
disjoint edges, each used twice,
\[
 |A_K|\le D_{N,h}\exp\left(\sum_{k=2}^K\Lambda_k\right),\qquad
 D_{n,k}=\frac{(n)_{2k}(2k)!}{4^k k!}.                  \tag{CB.3}
\]
\end{lemma}
\begin{proof}
Mark a component with $2k$ gates, retain its positions and all exterior
gates, and replace it by $k$ disjoint doubled edges using labels absent
from the exterior.  There are $n\ge N-2m$ available labels and precisely
$D_{n,k}$ replacements.  Each resulting word has at most $\binom hk$
reverse markings, since the reverse selects $k$ of its doubled-edge
components.  For a fixed reverse marking there are $U_{n,k}$ originals.
Giving each forward replacement weight $1/D_{n,k}$ handles the possible
variation of $n$.  Double counting gives
\[
 \mathbb E[Z_k\mid A_K]\le\binom hk
                  \sup_{N-2m\le n\le N}\frac{U_{n,k}}{D_{n,k}}.
\]
Take $c_0$ small enough that $n\ge N/2$ and $2k\le n/2$.
Use $(n)_{2k}\ge(n/2)^{2k}$,
$\binom hk k!\le h^k$, and
$\sum_{g=0}^k x^g\le2^k\max\{1,x\}^k$ to obtain (CB.2).

For the second assertion fix the vector $u=(Z_2,\ldots,Z_K)$ and denote
its class by $A_u$.  The same weighted switching gives
$u_k|A_u|\le\Lambda_k|A_{u-e_k}|$.  Iterating yields
$|A_u|\le D_{N,h}\prod_{k=2}^K\Lambda_k^{u_k}/u_k!$.
Summing over feasible vectors, then over all nonnegative vectors, proves
(CB.3).  The all-zero class is exactly the set of disjoint doubled edges.
\end{proof}

\subsection{Balancing the cut and bounding its flow}

\begin{proof}[Proof of Theorem~\ref{thm:balanced-bottleneck}]
Choose small fixed $a,c_0>0$ such that (CB.2), for
\[
 K_N=\min\{h,\lfloor aN/\sqrt m\rfloor\},
\]
gives $\Lambda_k\le C_0N4^{-k}$ for $2\le k\le K_N$ whenever
$m\le c_0N$.  Splitting the sum at $\sqrt N$ gives
\[
 \sum_{k=2}^{K_N}\Lambda_k
   =O(m^2/N)+O(N4^{-\sqrt N}).
\]
Indeed below $\sqrt N$ the ratio in (CB.2) is $C_1m/N$, and above it
the bound $C_0N4^{-k}$ applies.  Since $D_{N,h}\le N^m m^{m/2}$,
(CB.3) implies $|A_{K_N}|^{1/m}\le C N\sqrt m$.

For $s=\lfloor m/\log m\rfloor\le N$, let $a_g$ count length-$h$
words on $[s]$ with product $g$.  Reversing a word inverts its product,
so Cauchy--Schwarz gives
\[
 |\Omega(N,m)|\ge\sum_{g\in S_s}a_g^2
       \ge\binom s2^{\!m}/s!,\qquad
 |\Omega(N,m)|^{1/m}\ge c\frac{m^2}{(\log m)^2}.       \tag{CB.4}
\]
Consequently at
$n_- =\lfloor b_0m^{3/2}/(\log m)^2\rfloor$, for sufficiently small
fixed $b_0>0$, the capped event has mass at most $2^{-m}$.  At
$n_+=\lceil m^{3/2}/a\rceil$ its cap is $h$ and its mass is one.
For large $m$ the whole interval satisfies $m\le c_0N$.

Increase the alphabet from $N$ to $N+1$.  Conditioned on avoiding the new
label, the larger uniform law is exactly the smaller one.  Symmetry bounds
the probability of using that label by $2m/(N+1)$.  Also $K_{N+1}-K_N$
is either zero or one.  If it is one, the extra event at the new alphabet
has probability at most
\[
 \Pr(A_{K_N+1})\mathbb E[Z_{K_N+1}\mid A_{K_N+1}]
       \le C_0(N+1)4^{-(K_N+1)}.
\]
Thus the change in capped mass is $o(1)$ uniformly on $[n_-,n_+]$.
Monotonicity is not required: at the first crossing of $1/2$, the mass
is $1/2+o(1)$.  Fix this $N$ and its cap $K$.
Then $K\asymp N/\sqrt m\gtrsim m/(\log m)^2$.

In a pair update with nonidentity product, the two distinct factors have
the same component vertex sets and weights before and after the move.
Only a square refresh can change the profile.  Removing a square cannot
increase any component weight; inserting a new square joins at most two
remaining components and adds two gates.  If a word in $A_K$ exits, some
old component therefore has half-length at least $K/2$.  The case where
both joined pieces belong to the same old component satisfies the same
bound (and the deleted two gates must be subtracted if they came from it).
For $K\ge4$, (CB.2) gives
\[
 \frac{Q_P(A_K,A_K^c)}{\pi(A_K)}
 \le\sum_{k\ge\lceil K/2\rceil}^{K}\Lambda_k
 \le C N4^{-K/2}.                                    \tag{CB.5}
\]

For a move changing at most $b$ slots, contract the old components,
including isolates.  The at most $b$ inserted edges can connect at most
$b+1$ old components into one new component.  All unchanged gates in it
come from those components, and at most $b$ gates are new.  An exit
therefore requires an old half-length strictly greater than
$(K-b/2)/(b+1)$.  When this threshold is at least one, the same argument
bounds the escape probability by
\[
 C N4^{-s_b},\qquad
 s_b=\left\lfloor\frac{K-b/2}{b+1}\right\rfloor+1.    \tag{CB.6}
\]
For fixed $b$, the factor $CN$ is absorbed into $\exp(-c_bN/\sqrt m)$.
This argument only uses stationary uniform weights and the permitted
changes, not heat-bath transition probabilities or reversibility.

Finally stationarity gives equal flows across the cut in both directions.
For $X_0\sim\pi$, a union bound over $t$ steps bounds the probability
of crossing the cut by $2tQ_L(A_K,A_K^c)$.  Equivalently, the average over
individual starting words of their probability of leaving their initial
side is at most this quantity.  With the stated $t_m$, it tends to zero;
Markov's inequality makes the leaving probability $o(1)$ for a
$1-o(1)$ fraction of starting words.  Their initial side has stationary
mass $1/2+o(1)$, which gives the asserted TV bound.  Since $N$ is
polynomially bounded in $m$, the lower bound exceeds every polynomial in
$N,m$.
\end{proof}

The same proof permits growing update budgets
$b=o(m/(\log m)^3)$: (CB.6) is then superpolynomially small, so the
mixing time is superpolynomial along the same family.  For lazy reversible
kernels the indicator test function also gives an exponentially small
spectral-gap upper bound for each fixed $b$.

\subsection{Support-profile expansion at long lengths}

The preceding family has $m/N\to0$.  The conclusion changes for cuts
determined only by support profiles when $m$ is much larger.  A profile
records the labeled vertex set and gate count of every component,
including isolates.  This definition and the following proof do not
require a Hurwitz-orbit classification theorem.

\begin{theorem}[Long-length profile conductance]
\label{thm:long-profile-conductance}
For some absolute $C$ and all sufficiently large $N$, at every even
$m\ge CN\log N$ the pair kernel satisfies
\[
 \varphi:=\min_{\substack{0<\pi(A)\le1/2\\
                         A\text{ a union of profiles}}}
       \frac{Q_P(A,A^c)}{\pi(A)}=\Theta(N^{-3}).
\]
Let $\delta$ be the minimum conductance of the restrictions of $P$ to
individual profiles, with exits changed to holds and singleton
conductance set to one.  Then
\[
 \Phi(P)\ge\frac{\delta\varphi}{\delta+\varphi+2}
          \ge cN^{-3}\delta.
\]
The bound leaves $\delta$ unestimated; it is not a polynomial mixing
theorem for the full walk at $m=N^2$.
\end{theorem}

\begin{proof}
Write $d_v=\binom v2$ and let $\mu_{v,r}$ be the product distribution
of $r$ independent uniform transpositions on $v\ge2$ labels.
We first establish the parity-aware bound
\[
 \max_{\operatorname{sgn}(x)=(-1)^r}
 \left|\frac{v!}{2}\mu_{v,r}(x)-1\right|
 \le \exp\bigl(v^2e^{-r/(v-1)}\bigr)-1,\qquad r\ge1. \tag{CB.7}
\]
In the character expansion, the trivial and sign representations supply
the uniform parity-class law.  Bound every other character by its
dimension $f_\lambda$, and pair conjugate diagrams, whose transposition
ratios have opposite signs.  For a positive ratio with first row $v-j$,
the contents formula gives
\[
 r_\lambda=
 \frac{\sum_i\binom{\lambda_i}{2}-\sum_i\binom{\lambda'_i}{2}}{d_v}
 \le1-\frac j{v-1}.
\]
Indeed the first numerator sum is at most $v(v-j-1)/2$.
Choosing the entries below the first row and their tableau gives
$\sum_{\lambda:\lambda_1=v-j}f_\lambda^2
 \le\binom vj^2j!\le v^{2j}/j!$.
Summing over $j\ge1$ proves (CB.7); zero ratios contribute nothing
for $r\ge1$.  The $v=2$ parity law is exact.

Fix a sufficiently large absolute $D$.  If an even
$\ell\ge Dv\log N$, the first half of an unrestricted identity bridge
of length $\ell$ or $\ell-2$ has density at most two relative to
independent gates, by (CB.7) on the remaining half and the whole bridge.
For $r$ independent gates the probability that their graph is disconnected
is at most
\[
 \sum_{j=1}^{\lfloor v/2\rfloor}\binom vj
       \left(1-\frac{j(v-j)}{d_v}\right)^r.
\]
For $r\ge Dv\log N/3$ this is $O(N^{-15})$, uniformly in
$2\le v\le N$, after increasing $D$.  Both unrestricted bridges are
therefore connected with probability $1-O(N^{-15})$.
At a specified adjacent pair of the length-$\ell$ bridge, the square
probability is exactly
\[
 \frac{\mu_{v,\ell-2}(\id)}{d_v\mu_{v,\ell}(\id)}
          =\frac{1+o(1)}{d_v}.
\]
Conditional on a square, its deletion leaves a uniform unrestricted
length-$(\ell-2)$ identity bridge.  Conditioning the original word to be
connected thus gives probability at least $c/d_v$ of a square whose
deletion leaves all $v$ labels connected, and probability at most
$N^{-10}/d_v$ of a square whose deletion disconnects them.
In particular the full-support connected profile has mass greater than
$1/2$ when $m\ge2DN\log N$.

Conditioned on a profile, the component words are independent uniform
connected words and their colors have a uniform interleaving.  Call a
component $i$ dense if its gate count $\ell_i$ is at least
$Dv_i\log N$.  The dense components carry at least $m/2$ gates.
From a profile with at least two components, select moves that remove a
square from a dense source $i$ without disconnecting its labels, then
insert a square joining it to another dense component.  If there is only
one dense component, allow every other component, including isolates, as
target.  The merger is dense: in the latter case its source has at least
$m-DN\log N\ge DN\log N$ gates.

At a uniformly chosen adjacent pair, the probability that both colors
are $i$ is $\ell_i(\ell_i-1)/(m(m-1))$.  Their component indices are
consecutive, so the preceding deletion estimate applies.  There is at
least one eligible target label, and a square refresh joins $i$ to a
target with probability at least $v_i/d_N$.  The selected forward rate,
including laziness, is at least
\[
 \frac{c}{m^2N^2}\sum_{i\text{ dense}}\frac{\ell_i^2}{v_i}
 \ge\frac{c(\sum_{i\text{ dense}}\ell_i)^2}
            {m^2N^2\sum_{i\text{ dense}}v_i}
 \ge c'N^{-3}.
\]
Every reverse of a selected move deletes a square that disconnects a
dense component.  Conditional on any profile, its total reverse rate
is at most $N^{-10}$, by the deletion estimate and averaging over colors.
Any union $A$ of profiles of mass at most $1/2$ excludes the dominant
full-support profile.  Its selected outgoing flow is at least
$c'N^{-3}\pi(A)$.  The portion ending inside $A$ is, by reversibility,
at most the reverse flow originating in $A$, namely $N^{-10}\pi(A)$.
Their difference proves the lower bound for $\varphi$.

For the upper bound let $A$ be the event that label $N$ is unused.
It has positive mass less than $1/2$.  Its conditional law is the
unrestricted bridge on $N-1$ labels, in which an adjacent square occurs
with probability $(1+o(1))/d_{N-1}$.  Only a square refresh can introduce
$N$, doing so with probability $2/N$.  The escape rate is consequently
$(1+o(1))/(N d_{N-1})=\Theta(N^{-3})$.

For an arbitrary cut $A$ with $\pi(A)\le1/2$, round it to the union $B$
of profiles where its conditional mass exceeds $1/2$, and put
$u=\pi(A\mathbin{\triangle} B)$.  Restricted flow gives
$Q_P(A,A^c)\ge\delta u$.
Also $\min(\pi(B),1-\pi(B))\ge\pi(A)-u$, even if $\pi(B)>1/2$.
Stationarity and the symmetric difference bound therefore give
\[
 Q_P(A,A^c)\ge Q_P(B,B^c)-2u
              \ge\varphi\pi(A)-(\varphi+2)u.
\]
Eliminating $u$ proves the assertion (trivially if $\delta=0$).
\end{proof}

\subsection{Connectivity and a coarse all-length lower bound}

\begin{lemma}[All-endpoint pair connectivity]
\label{lem:all-endpoint-pair-connectivity}
For $N\ge2$, any two transposition words of the same length and product
are connected by adjacent same-product pair replacements.
\end{lemma}
\begin{proof}
In the active word, eliminate its largest available label $v$.  Carry
the first factor $t=(v,a)$ to the right.  On meeting $u$ fixing $v$, use
$(t,u)\mapsto(u,utu)$; on meeting $u=(v,b)\ne t$, use
$(t,u)\mapsto(tut,t)$.  In either case the left factor now fixes $v$.
On meeting another $t$, move the resulting square to reserved positions
at the front and continue.  The identity
$(t,t,u)\mapsto(t,tut,t)\mapsto(u,t,t)$ moves a square past a gate
without changing either one, using two inverse Hurwitz moves; reversing
this sequence moves it to the front.

If a carried gate reaches the end, the preceding active factors fix
$v$.  For active product $g$, the final gate must be
$(v,g^{-1}(v))$.  This happens exactly when $g$ moves $v$.  Freeze this
gate, update the active product to $g(v,g^{-1}(v))$, which fixes $v$,
and descend through the labels.  If $g$ fixes $v$, all its occurrences
are removed into reserved squares instead.  The frozen residual depends
only on the original endpoint.  The number of reserved squares depends
only on endpoint and length.  Recolor them all to $(1\,2)^2$ by pair
refreshes.  This is a common normal form; it proves connectivity only.
\end{proof}

\begin{proposition}[Coarse conductance bound]
\label{prop:coarse-pair-conductance}
For $N\ge3$ and every even $m\ge2$,
\[
 \Phi(P_{N,m})\ge\frac1{4N(m-1)d_N^{2N+1}}.           \tag{CB.8}
\]
\end{proposition}
\begin{proof}
Every nontrivial pair transition has probability at least
$p_*=1/[2(m-1)d_N]$: a same-product pair fiber has size $d_N$, three,
or two.  For $m>N$, put $W=N-1$.  For each ordered pair of identity
words $x,y$, construct splices with the first $i$ gates of $y$, a
length-$W$ connector, and the suffix of $x$ after position $i+W$.
The required connector permutation has parity $(-1)^W$ and minimum
transposition length at most $N-1$, so squares pad it to length $W$.
At $i=0$ and $i=m-W$ choose the connectors giving $x$ and $y$ themselves.
Consecutive splices differ in a window of $N$ gates, and the preceding
lemma joins them within it.  Fix all choices deterministically.

For a directed edge on such a route, there are at most $N$ possible
stage windows containing its two altered positions.  Fix one.  Outside
it the edge specifies the prefix of $y$ and suffix of $x$.  Encode the
complementary prefix of $x$ and suffix of $y$ in an identity word using
a connector of length $N$ (again parity permits it).  Recording also
the original length-$N$ windows of $x$ and $y$ recovers the ordered
pair.  Hence an edge can be used by at most
$N|\Omega(N,m)|d_N^{2N}$ ordered pairs.

Every pair directed from a cut to its complement uses an exiting edge.
For $|A|\le|\Omega|/2$, counting such pairs, then multiplying the
number of exiting edges by $p_*/|\Omega|$, gives (CB.8).
Only edge use, not path length or visit multiplicity, is counted here.
For $m\le N$, connectivity supplies an exiting edge for every nontrivial
cut and $|\Omega|\le d_N^m$ gives an even stronger bound.
\end{proof}

The bound is exponential in $N\log N$ and is consistent with the
bottleneck theorem.  Odd identity lengths are infeasible; $m=0$ and
$N=2$ give singleton identity spaces.  At $m=2$ the pair kernel refreshes
the whole doubled edge, whereas the project kernel has no permitted
block and must be treated as the identity kernel.  These exceptions are
outside Theorem~\ref{thm:balanced-bottleneck}.
